\documentclass[10pt,a4paper]{amsart}
\usepackage[margin=19mm]{geometry}
\usepackage{amsmath,amssymb,mathtools}
\usepackage[scr=rsfs,cal=euler]{mathalfa}
\usepackage{xfrac}
\usepackage[unicode,hidelinks,bookmarks=false]{hyperref}
\newcommand{\ie}{i.e.\ }
\newcommand{\cf}{cf.\ }
\newcommand{\resp}{respectively\ }
\newcommand{\Subsection}{\subsection}
\providecommand{\nobreakdash}{\nobreak\hbox{-}}
\setlength{\emergencystretch}{2em}
\multlinegap0pt
\usepackage{bm}
\usepackage{bbm}
\usepackage{dsfont}
\usepackage{xspace}
\usepackage{cancel}
\usepackage{mathtools}
\usepackage[shortlabels]{enumitem}
\usepackage{amsthm}
\makeatletter
\@ifundefined{equation}{\newtheorem{equation}{Equation}}{}
\@ifundefined{prop}{\newtheorem{prop}[equation]{Proposition}}{}
\makeatother
\newcommand{\C}{\mathbb{C}}          
\newcommand{\Hom}{\operatorname{Hom}}
\newcommand{\Z}{\mathbb{Z}}
\newcommand{\eps}{\varepsilon}
\newcommand{\fg}{{\mathfrak g}}
\newcommand{\fh}{{\mathfrak h}}
\newcommand{\fp}{{\mathfrak p}}
\newcommand{\ga}{\alpha}
\newcommand{\id}{\mathrm{id}}
\newcommand{\wlambda}{\widetilde{\lambda}}
\title[On the intertwining differential operators between vector bu]{On the intertwining differential operators between vector bundles over the real projective space of dimension two}
\author{Toshihisa Kubo, Bent {\O}rsted}
\date{Source: arXiv:2503.22323v3 (CC BY 4.0)}
\begin{document}
\maketitle

\noindent\emph{Source and licence.} On the intertwining differential operators between vector bundles over the real projective space of dimension two. Version arXiv:2503.22323v3 (CC BY 4.0), distributed under the Creative Commons Attribution 4.0 International licence, as recorded in the arXiv licence metadata for this exact version and shown on the abstract page https://arxiv.org/abs/2503.22323v3. Full rights notice: licenses/arxiv-2503.22323-CC-BY-4.0.txt.\par\smallskip

\noindent\textit{Editorial scope.} Counted unit: the Prop whose source label is \texttt{prop:Verma} in the article (source0.tex lines 1618--1627, source label \texttt{prop:Verma}), delivered together with the article's own complete original proof of it (source0.tex lines 1629--1794, one \texttt{proof} environment of 166 lines). No mathematics is written, repaired or re-derived: statement and proof are the article's own text line for line. Required context retained uncounted: eqn:Weyl (lines 1563--1570); eqn:HomVerma (lines 1603--1610); eqn:cond (lines 1612--1614). Every definition, display and label of those lines is retained unchanged. Omitted from this package and not redistributed: 1--1562, 1571--1602, 1611--1611, 1615--1617, 1628--1628, 1795--4030. Nothing inside a retained excerpt is omitted or altered: every retained body is delivered exactly as the article's lines are written, so no omission receipt of any kind accompanies this package. Citations: the retained excerpts carry no citation command at all, so no bibliography entry of the article is transcribed and no reference is redistributed. Reference adaptation: none was needed and none was made; every reference inside the delivered unit resolves inside it, so no unresolved reference or citation is produced. Printed slips kept literal: no printed word, symbol, hypothesis or number of the counted statement or of its proof was altered. Layout and class: the article's own class is \texttt{amsart} with packages including latexsym, float, amssymb, amsmath, amsthm, mathrsfs, euscript, xy, bbold, tabu, enumerate, framed, graphicx, stackrel; the wrapper is the lane's pinned \texttt{amsart} editorial wrapper at 10pt on A4, which restates the theorem-like environments the retained text needs and adds the article's own macro definitions that the retained text needs (\texttt{\textbackslash C}, \texttt{\textbackslash Hom}, \texttt{\textbackslash Z}, \texttt{\textbackslash eps}, \texttt{\textbackslash fg}, \texttt{\textbackslash fh}, \texttt{\textbackslash fp}, \texttt{\textbackslash ga}, \texttt{\textbackslash id}, \texttt{\textbackslash wlambda}), with extra wrapper packages (bm, bbm, dsfont, xspace, cancel, mathtools, enumitem). The delivered numbering is the unit's own. Assets: no retained span contains a figure environment, a graphics-inclusion command, a PDF-page inclusion, a table or a TikZ picture; no image file is copied into this package, the package declares no asset dependency and no image file is redistributed with it. Licence: version arXiv:2503.22323v3 of this article is distributed under the Creative Commons Attribution 4.0 International licence, as recorded in the arXiv licence metadata for this exact version and shown on the abstract page https://arxiv.org/abs/2503.22323v3. A full rights notice is delivered at licenses/arxiv-2503.22323-CC-BY-4.0.txt. Characters: every retained excerpt is pure ASCII, so the delivered engine, XeLaTeX with its default Latin Modern text font, reports no missing character.
\begin{equation}\label{eqn:Weyl}
W = \{e, 
s_{\eps_1-\eps_2}, \,
s_{\eps_2-\eps_3}, \,
s_{\eps_2-\eps_3}s_{\eps_1-\eps_2}, \,
s_{\eps_1-\eps_2}s_{\eps_2-\eps_3}, \,
s_{\eps_1-\eps_2}s_{\eps_2-\eps_3}s_{\eps_1-\eps_2}\}
\end{equation}

\begin{equation}\label{eqn:HomVerma}
\Hom_{\fg,P}
%\left(
(N_\fp(\ell\omega_1 -(\lambda +\tfrac{3}{2}k)d\chi+\rho)^{\ga+k},
N_\fp(m\omega_1 -\lambda d\chi+\rho)^\ga)
%\right) 
\neq \{0\}
\end{equation}

\begin{equation}\label{eqn:cond}
\frac{m+k-\ell}{2} \in \Z\cap [0,\min(k,m)].
\end{equation}

\begin{prop}\label{prop:Verma}
If \eqref{eqn:HomVerma} holds, then one of the following conditions holds.
\begin{enumerate}
\item[\emph{(a)}] $(k, \ell) = (0, m)$ and $\lambda \in \C$.
\item[\emph{(b)}] 
$\ell = m+k$ and $\lambda = \tfrac{1}{2}(2-m-2k)$.% (Cartan component)
\item[\emph{(c)}] 
$\ell=m-k$ and $\lambda = \tfrac{1}{2}(4+m-2k)$.% (PRV component)
\end{enumerate}
\end{prop}

\begin{proof}
By assumption,
the infinitesimal character of 
$N_\fp(\ell\omega_1 -(\lambda +\tfrac{3}{2}k)d\chi+\rho)^{\ga+k}$
agrees with that of $N_\fp(m\omega_1 -\lambda d\chi+\rho)^\ga$. 
Thus, there exists $w \in W$ such that
\begin{equation}\label{eqn:inf}
\ell\omega_1 -(\lambda +\tfrac{3}{2}k)d\chi+\rho
=w(m\omega_1 -\lambda d\chi+\rho).
\end{equation}

For a simple expression of both sides of \eqref{eqn:inf}, we normalize $\lambda$ by
\begin{equation*}
\wlambda:=\tfrac{2}{3}\lambda.
\end{equation*}
Then the weights $m\omega_1 -\lambda d\chi+\rho$ and 
$\ell\omega_1 -(\lambda +\tfrac{3}{2}k)d\chi+\rho$ are given in coordinates as
\begin{align*}
m\omega_1 -\lambda d\chi+\rho
&=m\omega_1 -\tfrac{3}{2}\wlambda d\chi+\rho\\
&=\tfrac{m}{2}(0, 1, -1) - \tfrac{\wlambda}{2}(2,-1,-1) + (1,0,-1)\\
&=
(1-\wlambda, \tfrac{1}{2}(\wlambda+m),\tfrac{1}{2}(\wlambda-m-2))
\end{align*}
and
\begin{align*}
\ell\omega_1 -(\lambda +\tfrac{3}{2}k)d\chi+\rho
&=\ell\omega_1 -\tfrac{3}{2}(\wlambda +k)d\chi+\rho\\
&=\tfrac{\ell}{2}(0,1,-1)-\tfrac{\wlambda+k}{2}(2,-1,-1)+(1,0,-1)\\
&=
(1-\wlambda-k, \tfrac{1}{2}(\wlambda+k+\ell), \tfrac{1}{2}(\wlambda+k-\ell-2)).
\end{align*}
Thus \eqref{eqn:inf}  amounts to
\begin{equation*}
(1-\wlambda-k, \tfrac{1}{2}(\wlambda+k+\ell), \tfrac{1}{2}(\wlambda+k-\ell-2))
=w(1-\wlambda, \tfrac{1}{2}(\wlambda+m),\tfrac{1}{2}(\wlambda-m-2)).
\end{equation*}

\noindent
Since the Weyl group element $s_{\eps_i-\eps_j}$ acts on $\fh^*\subset \C^3$
as the exchange of the $i$th entry with the $j$th entry, 
it follows from \eqref{eqn:Weyl} that one of the following cases holds.

\begin{enumerate}[(1)]

\item $1-\wlambda =1-\wlambda-k$.

\vskip 0.05in

\item $1-\wlambda = \tfrac{1}{2}(\wlambda+k+\ell)$.

\vskip 0.05in

\item $1-\wlambda = \tfrac{1}{2}(\wlambda+k-\ell-2)$.

\vskip 0.05in

\end{enumerate}
\noindent
We consider these cases separately.

\vskip 0.1in

%%%%%%%%%%%%%%%%%%%%%%%%%%%%%%%%%%%%%%%%%%
\noindent
Case (1): 
Since $1-\wlambda =1-\wlambda-k$,
we have $k=0$.
It then follows from \eqref{eqn:cond} that  $m=\ell$. Thus,
\begin{equation*}
N_\fp(\ell\omega_1 -(\lambda +\tfrac{3}{2}k)d\chi+\rho)^{\ga+k}=
N_\fp(m\omega_1 -\lambda d\chi+\rho)^\ga.
\end{equation*}
In this case, for any $\lambda\in \C$, we have
\begin{equation}\label{eqn:Case1}
%\Hom_{\fg}
%\left(
%N_\fp(\ell\omega_1 -(\lambda +\tfrac{3}{2}k)d\chi+\rho),
%N_\fp(m\omega_1 -\lambda d\chi+\rho)
%\right)
%&=
0\neq \id \in 
\Hom_{\fg, P}
%\left
(
N_\fp(m\omega_1 -\lambda d\chi+\rho)^{\ga},
N_\fp(m\omega_1 -\lambda d\chi+\rho)^\ga
%\right
),
\end{equation}
where $\id$ denotes the identity map.
\vskip 0.1in

%%%%%%%%%%%%%%%%%%%%%%%%%%%%%%%%%%%%%%%%%%
\noindent
Case (2): 
In this case, we have either
\begin{equation}\label{eqn:2a}
\begin{aligned}
1-\wlambda &= \tfrac{1}{2}(\wlambda+k+\ell),\\
 \tfrac{1}{2}(\wlambda+m) &=1-\wlambda-k,\\
\tfrac{1}{2}(\wlambda-m-2)&=\tfrac{1}{2}(\wlambda+k-\ell-2),
\end{aligned}
\end{equation}
or
\begin{equation}\label{eqn:2b}
\begin{aligned}
1-\wlambda &= \tfrac{1}{2}(\wlambda+k+\ell),\\
 \tfrac{1}{2}(\wlambda+m) &=\tfrac{1}{2}(\wlambda+k-\ell-2),\\
\tfrac{1}{2}(\wlambda-m-2)&=1-\wlambda-k.
\end{aligned}
\end{equation}

If \eqref{eqn:2a} holds, then a direct computation shows that
\begin{equation}\label{eqn:2a2}
\begin{aligned}
\wlambda&=\frac{1}{3}(2-m-2k),\\
\ell&=m+k.
\end{aligned}
\end{equation}

If \eqref{eqn:2b} holds, then $\ell=k-m-2$. 
Since $\ell \geq 0$, this implies that $\min(k,m)=m$.
It then follows from \eqref{eqn:cond} that there exists $d \in \{0,\ldots, m\}$ such that
$\ell = k+m-2d$. Thus,
\begin{equation*}
k-m-2=\ell =k+m-2d \geq k-m,
\end{equation*}
which is a contradiction. Therefore, \eqref{eqn:2b} does not occur. 

\vskip 0.1in

\item 
Case (3): 
As in Case (2), in this case, we have either
\begin{equation}\label{eqn:3a}
\begin{aligned}
1-\wlambda &= \tfrac{1}{2}(\wlambda+k-\ell-2),\\
\tfrac{1}{2}(\wlambda+m) &=\tfrac{1}{2}(\wlambda+k+\ell),\\
\tfrac{1}{2}(\wlambda-m-2)&=1-\wlambda-k,
\end{aligned}
\end{equation}
or
\begin{equation}\label{eqn:3b}
\begin{aligned}
1-\wlambda &= \tfrac{1}{2}(\wlambda+k-\ell-2),\\
\tfrac{1}{2}(\wlambda+m)&=1-\wlambda-k,\\
\tfrac{1}{2}(\wlambda-m-2) &=\tfrac{1}{2}(\wlambda+k+\ell).
\end{aligned}
\end{equation}

If \eqref{eqn:3a} holds, then we have
\begin{equation}\label{eqn:3a2}
\begin{aligned}
\wlambda&=\frac{1}{3}(4+m-2k),\\
\ell&=m-k.
\end{aligned}
\end{equation}

If \eqref{eqn:3b} holds, then $\ell=-(m+k+2) < 0$, which contradicts 
the condition that  $\ell \geq 0$. 
Thus, this case does not happen.

The conditions \eqref{eqn:Case1}, \eqref{eqn:2a2},
and \eqref{eqn:3a2} conclude the proposition.
\end{proof}

\end{document}
